\chapter{Conclusion et perspectives}

\section{Bilan}

L'application couvre l'ensemble du cahier des charges : les 23 besoins fonctionnels (P1, P2 et P3), les
16 règles de gestion, contrôlées dans la couche service, et les besoins non fonctionnels de sécurité, de
fiabilité et d'ergonomie. La démarche du modèle en Y a permis de fixer les cas d'utilisation et la
conception avant le code, puis de faire évoluer les diagrammes de conception avec l'implémentation.
Le passage de commande, transactionnel et protégé contre les accès concurrents, constitue le cœur
technique du projet.

\section{Perspectives}

Améliorations possibles, hors du périmètre actuel :

\begin{itemize}
  \item Paiement en ligne réel (le paiement est aujourd'hui simulé, à la livraison).
  \item Envoi d'e-mails transactionnels (confirmation de commande, changement de statut).
  \item Suivi de livraison avec un transporteur.
  \item Panier persistant en base, conservé entre deux connexions.
  \item Envoi d'images de produits (aujourd'hui une URL d'image est saisie).
  \item Contrôle de l'état du compte à chaque requête, pour déconnecter immédiatement un client désactivé.
  \item Signalement dans le panier des produits devenus indisponibles avant la validation.
  \item Tests automatisés (tests unitaires des services, tests d'intégration des DAO).
  \item Pool de connexions de production (par exemple HikariCP) à la place du pool intégré d'Hibernate.
\end{itemize}
